\documentclass[aspectratio=169,10pt]{beamer}
\usetheme[numbering=fraction,progressbar=frametitle]{metropolis}
\usepackage{fontspec}
\setsansfont{Latin Modern Sans}
\usepackage{amsmath,amssymb,mathtools,graphicx,booktabs}
\usepackage{polyglossia}
\setdefaultlanguage{english}
\setbeamertemplate{navigation symbols}{}
\setbeamerfont{frametitle}{size=\large}
\setbeamerfont{normal text}{size=\normalsize}
\setlength{\parskip}{.45em}
\hypersetup{pdfauthor={},pdfsubject={Lecture slides},pdfkeywords={},pdftitle={Random variables and dispersion}}
\providecommand{\N}{\mathbb{N}}
\providecommand{\Z}{\mathbb{Z}}
\providecommand{\D}{\mathbb{D}}
\providecommand{\Q}{\mathbb{Q}}
\providecommand{\R}{\mathbb{R}}
\providecommand{\C}{\mathbb{C}}
\providecommand{\vect}[1]{\overrightarrow{#1}}
\providecommand{\vi}{\vec{\imath}}
\providecommand{\vj}{\vec{\jmath}}
\providecommand{\vk}{\vec{k}}
\providecommand{\norm}[1]{\left\lVert #1\right\rVert}
\providecommand{\vc}[2]{\begin{pmatrix}#1\\#2\end{pmatrix}}
\providecommand{\deter}[2]{\det\left(#1,#2\right)}
\providecommand{\abs}[1]{\left\lvert #1\right\rvert}
\providecommand{\dd}{\,\mathrm{d}}
\providecommand{\e}{\mathrm{e}}
\providecommand{\Card}{\operatorname{Card}}
\providecommand{\E}{\operatorname{E}}
\providecommand{\V}{\operatorname{V}}
\providecommand{\degre}{^{\circ}}
\providecommand{\un}[1]{\,\mathrm{#1}}
\providecommand{\sci}[2]{#1\times10^{#2}}
\providecommand{\pgcd}{\operatorname{PGCD}}
\providecommand{\ppcm}{\operatorname{PPCM}}
\title{Random variables and dispersion}
\subtitle{Lecture and worked examples}
\date{}
\author{}
\begin{document}
\begin{frame}[plain]
\vfill
{\usebeamerfont{title}\inserttitle\par}
\medskip
{\usebeamerfont{subtitle}\insertsubtitle\par}
\medskip
{\usebeamercolor[fg]{progress bar}\rule{\textwidth}{.5pt}}
\vfill
\end{frame}
\begin{frame}{Chapter path}
\textbf{Objectives.} Build a distribution; interpret events; calculate expectation, variance, standard deviation; use linearity; simulate samples and calculate their observed mean.

\textbf{Prerequisites.} Probability, weighted mean, squares, trees and lists.

\textbf{Yearly reference point.} 19-28 January: 8 hours
\end{frame}
\begin{frame}{Variable, event and distribution}
On a finite sample space \(\Omega\), a real random variable X is a function that associates a real number with each outcome. The event \(\{X=a\}\) groups the outcomes giving the value a. Its distribution associates with each distinct value \(x_i\) the probability \(p_i=P(X=x_i)\), with \(p_i\geq0\) and \(\sum p_i=1\).
\end{frame}
\begin{frame}{Variable, event and distribution}
\begin{columns}[T,onlytextwidth]
\begin{column}{.57\textwidth}
A fair die gives a gain X equal to 2 for 1 or 2, and to \(-1\) otherwise. The distribution is \(P(X=2)=\dfrac{1}{3}\), \(P(X=-1)=\dfrac{2}{3}\). \(\{X\leq0\}=\{3,4,5,6\}\): probability \(\dfrac{2}{3}\).
\end{column}
\begin{column}{.4\textwidth}
\centering
\includegraphics[width=\linewidth,height=4.5cm,keepaspectratio]{figures/premiere-08-1-en.pdf}
\end{column}
\end{columns}
\end{frame}
\begin{frame}{Variable, event and distribution}
The values are not always equally likely even if the outcomes are. Group the outcomes that give the same gain.
\end{frame}
\begin{frame}{Variable, event and distribution}
Two fair coins are tossed. X counts the heads. Write its distribution and \(P(X\leq1)\).
\end{frame}
\begin{frame}{Expectation and linearity}
The expectation is \(E(X)=\sum_i x_ip_i\). It describes the long-term theoretical mean, not a result that is necessarily possible on a single trial. For variables X and Y on the same finite sample space and real numbers a,b: \(E(aX+b)=aE(X)+b\) and \(E(X+Y)=E(X)+E(Y)\), with no assumption of independence.
\end{frame}
\begin{frame}{Expectation and linearity}
For the previous distribution \(P(X=2)=\dfrac{1}{3}\), \(P(X=-1)=\dfrac{2}{3}\), \(E(X)=0\): the game is fair if X is the net gain. For \(Y=3X-2\), \(E(Y)=3\times0-2=-2\).
\end{frame}
\begin{frame}{Expectation and linearity}
An expectation is not a guarantee for an individual: the accumulated outcome can be a loss at every observed stage even when the expectation is positive.
\end{frame}
\begin{frame}{Expectation and linearity}
A gross gain is 10 with probability \(\dfrac{1}{5}\) and 0 otherwise. What stake makes the game fair?
\end{frame}
\begin{frame}{Variance, standard deviation and the König-Huygens formula}
Writing \(\mu=E(X)\), \(V(X)=E((X-\mu)^2)=\sum_i(x_i-\mu)^2p_i\). The standard deviation is \(\sigma(X)=\sqrt{V(X)}\). The König-Huygens formula gives \(V(X)=E(X^2)-E(X)^2\). A variance is always nonnegative; the standard deviation has the same unit as X, while the variance has the square of that unit.
\end{frame}
\begin{frame}{Variance, standard deviation and the König-Huygens formula}
For X with values 0 and 4 with probabilities \(\dfrac{3}{4}\) and \(\dfrac{1}{4}\), \(E(X)=1\), \(E(X^2)=4\), \(V(X)=3\) and \(\sigma=\sqrt3\). Direct calculation: \((0-1)^2(\dfrac{3}{4})+(4-1)^2(\dfrac{1}{4})=3\).
\end{frame}
\begin{frame}{Variance, standard deviation and the König-Huygens formula}
\(E(X^2)\) is not \(E(X)^2\). A variance calculated as negative signals an error or a very small rounding error to be dealt with explicitly.
\end{frame}
\begin{frame}{Variance, standard deviation and the König-Huygens formula}
For X with values \(-1,1\) equally likely, calculate the expectation, variance and standard deviation.
\end{frame}
\begin{frame}{Transforming and comparing gains}
For real a,b, \(V(aX+b)=a^2V(X)\) and \(\sigma(aX+b)=|a|\sigma(X)\). A translation moves the mean without changing the dispersion; a multiplication changes the scale. Two games with the same expectation can have very different dispersions.
\end{frame}
\begin{frame}{Transforming and comparing gains}
A game A always pays 2. A game B pays 0 or 4 with equal probability. Both have expectation 2; variances 0 and 4. Their distributions are therefore not interchangeable despite the same mean.
\end{frame}
\begin{frame}{Transforming and comparing gains}
The variance of a sum is not, in general, the sum of the variances; no rule of this type is used without hypotheses.
\end{frame}
\begin{frame}{Transforming and comparing gains}
\(E(X)=3\), \(V(X)=4\). Give the characteristics of \(Y=-2X+5\).
\end{frame}
\begin{frame}{Samples and observed mean}
A sample of size n is here a list of n independent realisations of the same variable X. Its observed mean m varies from one sample to another; we compare it with \(\mu=E(X)\).
\end{frame}
\begin{frame}{Samples and observed mean}
For X taking the values 0 and 2 with equal probability, \(\mu=1\). The sample 2, 0, 2, 2, 0, 0, 2, 2, 0, 2 has observed mean 1.2; another sample of the same size will usually give another mean. A simulation compares several observed means with \(\mu\) without proving anything.
\end{frame}
\begin{frame}{Samples and observed mean}
Do not copy a single draw n times; do not confuse the observed mean of a sample with the expectation \(\mu\).
\end{frame}
\begin{frame}{Samples and observed mean}
\begin{columns}[T,onlytextwidth]
\begin{column}{.57\textwidth}
X takes the values 0 and 2 with equal probability. Calculate the observed mean of the sample 2, 0, 2, 2, 0, then its gap from \(E(X)\).
\end{column}
\begin{column}{.4\textwidth}
\centering
\includegraphics[width=\linewidth,height=4.5cm,keepaspectratio]{figures/premiere-08-2-en.pdf}
\end{column}
\end{columns}
\end{frame}
\end{document}
